\section{Conclusion}
\label{sec:conclusion}

Court orders are a powerful instrument for securing access to
medicines. They also arrive as procurement shocks. In S\~ao Paulo,
those shocks raise prices and reduce bidder participation, even as
the state completes urgent tenders more often. The administrative
urgent channel shows that the court-sanction margin remains positive
after selection bounds, but the mechanism is not a broad incumbent
markup in deep repeated urgent markets. The main loss is sourcing
efficiency: judicial urgency fragments demand, reduces scale, and
changes which suppliers win.

That distinction matters for policy. If the procurement-cost margin
were mainly same-firm pricing, price caps or contract terms would be
the natural tools. If the margin is fragmented sourcing, the relevant
tools are those that preserve aggregation while respecting judicial
delivery obligations: framework agreements, pooled urgent
procurement, pre-contracted suppliers, inventories for recurrent
medicines, and administrative screening capacity. The point is not to
weaken access, but to prevent legal urgency from forcing the state to
buy one patient-specific emergency at a time.

The estimates are specific to S\~ao Paulo pharmaceutical procurement
during \BPsampleStartYear--\BPsampleEndYear{} and do not value health
benefits or the full social value of litigation. They identify a narrower but policy-relevant
margin: how judicial enforcement changes the way the state buys.
For public economics, that margin is central. Legal rights are
implemented through administrative production systems. When
enforcement changes those systems' timing, scale, and supplier
matches, the fiscal cost of a right depends not only on what is
mandated, but on how the state must procure it.
